\documentclass[12pt]{article}
\usepackage[utf8]{inputenc}
\usepackage[spanish]{babel}
\usepackage{amsmath}
\usepackage{amsfonts}
\usepackage{amssymb}
\usepackage{geometry}
\geometry{letterpaper, margin=2.5cm}

\begin{document}

\begin{center}
    \Large \textbf{Evaluación de Examen 1: Unidad I} \\
    \large Física para Ciencias de la Salud (005-1713) \\
    \vspace{0.5cm}
    \textbf{Estudiante:} Camila Cabrera \quad \textbf{C.I.:} 32.937.064
    \vspace{0.3cm} \\
    \large \textbf{Calificación Final: 9/10 puntos}
\end{center}

\vspace{0.5cm}

\section*{Análisis de Resultados}

\subsection*{1. Ángulo plano (Conversión a radianes)}
\textbf{Procedimiento y resultado:} Plantea una regla de tres a partir de la equivalencia $180^\circ = \pi \text{ rad}$. Realiza el desarrollo aritmético multiplicando por el valor decimal de $\pi$ ($3,14159$) para obtener un resultado exacto en notación decimal de $1,30 \text{ rad}$. Procedimiento lógico y correcto. \\
\textbf{Veredicto:} \textbf{Correcto} (2/2 pts).

\subsection*{2. Sistemas de coordenadas (Longitud del hueso)}
\textbf{Procedimiento y resultado:} Utiliza la fórmula formal de la distancia euclidiana entre dos puntos. Desarrolla las operaciones paso a paso ($\sqrt{6^2 + 8^2} = \sqrt{100}$) llegando sin errores al resultado de $10 \text{ cm}$. Además, la construcción de la gráfica y la ubicación del segmento en el plano son impecables. \\
\textbf{Veredicto:} \textbf{Correcto} (2/2 pts).

\subsection*{3. Vectores (Fuerza resultante)}
\textbf{Procedimiento y resultado:} Demuestra mucho orden al calcular por separado la suma de las componentes en el eje $x$ ($45 - 15 = 30$) y en el eje $y$ ($25 + 35 = 60$). Esto le permite armar correctamente el vector de fuerza resultante final: $\vec{F}_R = 30\hat{i} + 60\hat{j}$ N. \\
\textbf{Veredicto:} \textbf{Correcto} (2/2 pts).

\subsection*{4. Potencias de diez (Notación científica y prefijo)}
\textbf{Procedimiento y resultado:} Transforma la cifra correctamente a notación científica estándar ($1,25 \times 10^{-5}$ m). Sin embargo, la estudiante no completó la segunda mitad de la instrucción que consistía en identificar e indicar el prefijo adecuado del Sistema Internacional (micrómetros). \\
\textbf{Veredicto:} \textbf{Incompleto / Parcialmente correcto} (1/2 pts).

\subsection*{5. Sistemas de unidades (Conversión de densidad)}
\textbf{Procedimiento y resultado:} Utiliza de forma magistral los factores de conversión en cadena $\left(\frac{1 \text{ kg}}{1000 \text{ g}}\right)$ y $\left(\frac{1.000.000 \text{ cm}^3}{1 \text{ m}^3}\right)$. Desarrolla la operación de las fracciones y simplifica claramente la división para obtener el resultado exacto de $1030 \text{ kg/m}^3$. \\
\textbf{Veredicto:} \textbf{Correcto} (2/2 pts).

\vspace{0.5cm}
\section*{Comentarios Generales y Sugerencias}
¡Excelente examen! La estudiante demuestra un gran manejo de los procedimientos matemáticos, siendo muy metódica y ordenada al desglosar sus operaciones (especialmente evidente en el cálculo de vectores y en la gráfica de coordenadas).
\begin{itemize}
    \item Se recomienda prestar un poco más de atención a la lectura de los enunciados para asegurar responder el 100\% de lo que se pide, tal como ocurrió al omitir el prefijo del S.I. en la Pregunta 4.
\end{itemize}

\end{document}